% ============================================================================
%  main.tex — Universal LaTeX Template (Entry Point)
%  Author:  Yazid TAHIRI ALAOUI
%  Usage:   Edit the metadata commands below, then compile with pdflatex.
% ============================================================================
\documentclass[a4paper,12pt]{report}

% ── Load the style package ──────────────────────────────────────────────────
\usepackage{styles/yazid-template}

% ============================================================================
%  DOCUMENT METADATA  — Edit these for each new project
% ============================================================================
\newcommand{\docKingdom}{ROYAUME DU MAROC}
\newcommand{\docUniversity}{Université Abdelmalek Essaâdi}
\newcommand{\docSchool}{École Nationale des Sciences Appliquées d'Al Hoceima}
\newcommand{\docTitle}{Rapport de TP N\textdegree~5}
\newcommand{\docSubtitle}{La Généricité \& La Gestion des Exceptions}
\newcommand{\docModule}{Programmation Orientée Objet -- Java SE}
\newcommand{\docAuthor}{Yazid TAHIRI ALAOUI}
\newcommand{\docProgram}{1\`{e}re année -- Ingénierie Données (ID1)}
\newcommand{\docSupervisor}{Pr. Mohamed CHERRADI}
\newcommand{\docYear}{2025 -- 2026}

% ── Headers (uses commands from the .sty) ───────────────────────────────────
\renewcommand{\headerLeft}{TP5 -- Généricité \& Exceptions}
\renewcommand{\headerRight}{ENSAH -- Année 2025/2026}

% ── PDF metadata ────────────────────────────────────────────────────────────
\hypersetup{
    pdftitle={Rapport TP5 -- Généricité et Gestion des Exceptions},
    pdfauthor={Yazid TAHIRI ALAOUI},
    pdfsubject={Programmation Orientée Objet Java SE},
    pdfkeywords={Java, Exceptions, Généricité, ENSAH}
}

% ============================================================================
%  BEGIN DOCUMENT
% ============================================================================
\begin{document}

% ── Title Page ──────────────────────────────────────────────────────────────
\input{frontmatter/titlepage}

% ── Front Matter (roman page numbers) ──────────────────────────────────────
\pagenumbering{roman}
\newgeometry{left=2.5cm, right=2.5cm, top=2.5cm, bottom=2.5cm, headheight=15pt}

% ── Abstract / Résumé ──────────────────────────────────────────────────────
\renewcommand{\thechapter}{\Roman{chapter}}
% ============================================================================
%  abstract.tex — Résumé / Abstract Template
%  Replace the placeholder text below with your project's summary.
% ============================================================================

% --- French Abstract ---
\chapter*{Résumé}
\addcontentsline{toc}{chapter}{Résumé}

\begin{center}
{\large
Ce rapport présente les travaux réalisés dans le cadre du TP N°5 du module
\textbf{Programmation Orientée Objet en Java SE}, portant sur deux concepts
fondamentaux du langage Java~: la \textbf{gestion des exceptions} et la
\textbf{généricité}.\\[0.4cm]

Le TP est organisé en trois parties~: la première traite des exceptions
système (vérifiées et non vérifiées) à travers des exemples concrets de
division, de saisie utilisateur, de tableaux et de propagation d'exceptions.
La deuxième partie porte sur la création d'exceptions personnalisées
appliquées à des cas métier tels que les comptes bancaires, la gestion de
stock et les formulaires. La troisième partie explore la généricité avec
des classes, interfaces, méthodes génériques, bornes et wildcards.\\[0.4cm]

\textbf{Mots-clés~:} Java, Exceptions, Généricité, Checked, Unchecked,
Wildcards, Programmation Orientée Objet.
\vfill
}
\end{center}

% --- English Abstract (optional, uncomment if needed) ---
% \chapter*{Abstract}
% \addcontentsline{toc}{chapter}{Abstract}
%
% \begin{center}
% {\large
% Provide a concise summary of your project in English here.\\[0.4cm]
%
% State the problem, methodology, and key findings.\\[0.4cm]
%
% \textbf{Keywords:} Keyword~1, Keyword~2, Keyword~3, Keyword~4.
% \vfill
% }
% \end{center}


% ── Table of Contents ──────────────────────────────────────────────────────
\newpage
\renewcommand{\contentsname}{Table des matières}
\tableofcontents
\newpage

% ── Main Content (arabic page numbers) ─────────────────────────────────────
\pagenumbering{arabic}
\setcounter{page}{1}
\renewcommand{\thechapter}{\arabic{chapter}}
\setcounter{chapter}{0}

% ── Include your chapters here ─────────────────────────────────────────────
\chapter{Exceptions Système}
\sloppy

\section{Exercice 1 -- Division Sécurisée}

\subsection{Objectif}

Calculer la division de deux entiers saisis par l'utilisateur en gérant
la division par zéro et la saisie invalide.

\subsection{Code source}

\begin{lstlisting}[language=Java, caption={exo1.java -- Division sécurisée}]
import java.util.Scanner;
import java.util.InputMismatchException;

public class exo1 {

    public exo1() {}

    public void calculer(exo1 e, int a, int b) {
        try {
            if (b == 0) {
                throw new ArithmeticException();
            }
            System.out.println(a / b);
        } catch (ArithmeticException ex) {
            System.out.println("Erreur : Division par zero.");
        } catch (Exception ex) {
            System.out.println("Erreur inattendue.");
        }
    }

    public static void main(String[] args) {
        Scanner sc = new Scanner(System.in);
        try {
            System.out.print("Enter a : ");
            int a = sc.nextInt();
            System.out.print("Enter b : ");
            int b = sc.nextInt();
            exo1 d = new exo1();
            d.calculer(d, a, b);
        } catch (InputMismatchException ex) {
            System.out.println("Erreur : Saisie invalide, entrez des entiers.");
        }
        sc.close();
    }
}
\end{lstlisting}

\subsection{Explication}

Le bloc \texttt{try/catch} dans la méthode \texttt{calculer()} intercepte
l'\texttt{ArithmeticException} levée explicitement via \texttt{throw} lorsque
le diviseur est nul. Dans le \texttt{main}, un second \texttt{try/catch}
capture l'\texttt{InputMismatchException} levée par \texttt{Scanner} en cas de
saisie non entière.

\begin{infobox}
    \centering\vspace{0.2cm}
    {\large Le mot-clé \texttt{throw} permet de lever manuellement une exception
    depuis n'importe quel point du code.}
    \vspace{0.2cm}
\end{infobox}

\section{Exercice 2 -- NullPointerException}

\subsection{Objectif}

Afficher la longueur d'une chaîne en gérant le cas où elle est \texttt{null},
d'abord sans \texttt{try/catch} puis avec.

\subsection{Code source}

\begin{lstlisting}[language=Java, caption={exo2.java -- NullPointerException}]
public class exo2 {

    public void longueurSansCatch(exo2 e, String s) {
        if (s == null) {
            System.out.println("Null");
        } else {
            System.out.println(s.length());
        }
    }

    public void longueurAvecCatch(exo2 e, String s) {
        try {
            System.out.println(s.length());
        } catch (NullPointerException ex) {
            System.out.println("Null");
        }
    }
}
\end{lstlisting}

\subsection{Explication}

La version sans \texttt{try/catch} utilise une condition \texttt{if} pour
tester la nullité avant d'appeler la méthode. La version avec \texttt{try/catch}
laisse la \texttt{NullPointerException} se lever et la capture dans le bloc
\texttt{catch}.

\section{Exercice 3 -- Tableau Sécurisé}

\subsection{Objectif}

Permettre à l'utilisateur d'accéder à un index d'un tableau de taille 5
en proposant deux approches.

\subsection{Code source}

\begin{lstlisting}[language=Java, caption={exo3.java -- Accès sécurisé au tableau}]
public class exo3 {

    public void versionIf(exo3 e, int[] tab, int index) {
        if (index >= 0 && index < 5) {
            System.out.println(tab[index]);
        } else {
            System.out.println("Hors limites");
        }
    }

    public void versionCatch(exo3 e, int[] tab, int index) {
        try {
            System.out.println(tab[index]);
        } catch (ArrayIndexOutOfBoundsException ex) {
            System.out.println("Hors limites");
        }
    }
}
\end{lstlisting}

\subsection{Comparaison des deux approches}

\begin{table}[H]
\centering
\begin{tabularx}{0.85\textwidth}{C{3cm} X X}
\toprule
\textbf{\color{primarycolor}Critère} &
\textbf{\color{primarycolor}Version if} &
\textbf{\color{primarycolor}Version try/catch} \\
\midrule
Lisibilité & Explicite & Nécessite de connaître l'exception \\
\addlinespace[0.15cm]
Performance & Légèrement plus rapide & Surcoût si exception levée \\
\addlinespace[0.15cm]
Usage recommandé & Vérification préventive & Cas imprévisibles \\
\bottomrule
\end{tabularx}
\caption{Comparaison entre vérification par \texttt{if} et \texttt{try/catch}}
\end{table}

\section{Exercice 4 -- Conversion d'Entier}

\subsection{Objectif}

Convertir une chaîne saisie en entier et afficher un message clair en cas d'échec.

\subsection{Code source}

\begin{lstlisting}[language=Java, caption={exo4.java -- NumberFormatException}]
public class exo4 {

    public void convertir(exo4 e, String s) {
        try {
            int val = Integer.parseInt(s);
            System.out.println(val);
        } catch (NumberFormatException ex) {
            System.out.println("Erreur : La chaine saisie ne peut pas etre convertie en entier.");
        }
    }
}
\end{lstlisting}

\subsection{Explication}

\texttt{Integer.parseInt()} lève une \texttt{NumberFormatException} si la
chaîne ne représente pas un entier valide. Ce type d'exception est une
\textit{unchecked exception} (hérite de \texttt{RuntimeException}).

\section{Exercice 5 -- Choix d'Exception}

\subsection{Objectif}

Créer une méthode \texttt{racineCarree(int x)} qui lève une exception si \texttt{x < 0}.

\subsection{Code source}

\begin{lstlisting}[language=Java, caption={exo5.java -- IllegalArgumentException}]
public class exo5 {

    public int racineCarree(exo5 e, int x) {
        if (x < 0) {
            throw new IllegalArgumentException();
        }
        return (int) Math.sqrt(x);
    }
}
\end{lstlisting}

\subsection{Justification du choix}

\begin{accentbox}
    \centering\vspace{0.3cm}
    {\large \texttt{IllegalArgumentException} est choisie car l'erreur provient
    d'un \textbf{argument invalide} passé par le développeur.
    C'est une \textit{unchecked exception} qui ne nécessite pas de
    \texttt{throws} dans la signature, ce qui est adapté à ce type de
    précondition logique.}
    \vspace{0.3cm}
\end{accentbox}

\section{Exercice 6 -- IllegalStateException}

\subsection{Objectif}

Simuler une machine avec un état ON/OFF~; lever une exception si on tente de
la démarrer alors qu'elle est déjà allumée.

\subsection{Code source}

\begin{lstlisting}[language=Java, caption={exo6.java -- IllegalStateException}]
public class exo6 {

    boolean estAllumee;

    public exo6() {
        this.estAllumee = false;
    }

    public void demarrer(exo6 e) {
        if (e.estAllumee) {
            throw new IllegalStateException();
        }
        e.estAllumee = true;
        System.out.println("Machine ON");
    }
}
\end{lstlisting}

\subsection{Explication}

\texttt{IllegalStateException} est l'exception standard Java à utiliser
lorsqu'une méthode est appelée à un moment où l'\textbf{état de l'objet}
ne le permet pas.

\section{Exercice 7 -- Propagation d'Exception}

\subsection{Objectif}

Illustrer la propagation d'exceptions à travers une chaîne d'appels
A $\rightarrow$ B $\rightarrow$ C.

\subsection{Code source}

\begin{lstlisting}[language=Java, caption={exo7.java -- Propagation}]
public class exo7 {

    public void methodeC(exo7 e) throws Exception {
        throw new Exception();
    }

    public void methodeB(exo7 e) throws Exception {
        e.methodeC(e);
    }

    public void methodeA(exo7 e) {
        try {
            e.methodeB(e);
        } catch (Exception ex) {
            System.out.println("Erreur interceptee par A");
        }
    }
}
\end{lstlisting}

\subsection{Schéma de propagation}

\begin{infobox}
    \centering\vspace{0.2cm}
    {\large \texttt{C} lance l'exception $\rightarrow$ \texttt{B} la propage
    (\texttt{throws}) $\rightarrow$ \texttt{A} la capture (\texttt{catch})}
    \vspace{0.2cm}
\end{infobox}

\section{Exercice 8 -- Checked Exception}

\subsection{Objectif}

Créer une méthode \texttt{verifierAge(int age) throws Exception} qui lève une
exception si l'âge est inférieur à 18.

\subsection{Code source}

\begin{lstlisting}[language=Java, caption={exo8.java -- Checked Exception}]
public class exo8 {

    public void verifierAge(exo8 e, int age) throws Exception {
        if (age < 18) {
            throw new Exception();
        }
    }

    public static void main(String[] args) {
        exo8 e = new exo8();
        try {
            e.verifierAge(e, 15);
            System.out.println("Age valide.");
        } catch (Exception ex) {
            System.out.println("Erreur : Age inferieur a 18.");
        }
    }
}
\end{lstlisting}

\subsection{Explication}

En utilisant \texttt{throws Exception}, le compilateur \textbf{oblige}
l'appelant à gérer l'exception. C'est le comportement des
\textit{checked exceptions}.

\section{Exercice 9 -- Exception Personnalisée Vérifiée}

\subsection{Objectif}

Créer une \texttt{AgeException} personnalisée et l'utiliser pour vérifier un âge.

\subsection{Code source}

\begin{lstlisting}[language=Java, caption={exo9.java -- AgeException personnalisée}]
class AgeException extends Exception {}

public class exo9 {

    public void verifierAge(exo9 e, int age) throws AgeException {
        if (age < 18) {
            throw new AgeException();
        }
    }

    public static void main(String[] args) {
        exo9 e = new exo9();
        try {
            e.verifierAge(e, 15);
            System.out.println("Age valide.");
        } catch (AgeException ex) {
            System.out.println("Erreur : Age inferieur a 18.");
        }
    }
}
\end{lstlisting}

\section{Exercice 10 -- Comparaison Checked / Unchecked}

\subsection{Objectif}

Comparer les deux approches~: \texttt{Exception} (checked) et \texttt{RuntimeException} (unchecked).

\subsection{Code source}

\begin{lstlisting}[language=Java, caption={exo10.java -- Comparaison Checked vs Unchecked}]
public class exo10 {

    public void verifierChecked(exo10 e, int age) throws Exception {
        if (age < 18) {
            throw new Exception();
        }
    }

    public void verifierUnchecked(exo10 e, int age) {
        if (age < 18) {
            throw new RuntimeException();
        }
    }
}
\end{lstlisting}

\subsection{Comparaison}

\begin{table}[H]
\centering
\begin{tabularx}{0.9\textwidth}{C{3.5cm} X X}
\toprule
\textbf{\color{primarycolor}Critère} &
\textbf{\color{primarycolor}Exception (Checked)} &
\textbf{\color{primarycolor}RuntimeException (Unchecked)} \\
\midrule
Vérification & À la compilation & Seulement à l'exécution \\
\addlinespace[0.15cm]
Signature & Oblige \texttt{throws} & Aucun \texttt{throws} requis \\
\addlinespace[0.15cm]
Usage & Erreurs récupérables externes & Bugs logiques du développeur \\
\bottomrule
\end{tabularx}
\caption{Comparaison entre \texttt{Exception} et \texttt{RuntimeException}}
\end{table}

\chapter{Exceptions Personnalisées}
\sloppy

\section{Exercice 1 -- Compte Bancaire}

\subsection{Objectif}

Créer une classe \texttt{CompteBancaire} avec \texttt{retirer()} et \texttt{verser()},
et lever une \texttt{SoldeInsuffisantException} si le montant dépasse le solde.

\subsection{Code source}

\begin{lstlisting}[language=Java, caption={exo1\_2.java -- SoldeInsuffisantException}]
class SoldeInsuffisantException extends Exception {}

public class exo1_2 {

    String code;
    double solde;

    public exo1_2(String code, double solde) {
        this.code = code;
        this.solde = solde;
    }

    void verser(exo1_2 e, double montant) {
        e.solde += montant;
        System.out.println("Nouveau solde : " + e.solde);
    }

    void retirer(exo1_2 e, double montant) throws SoldeInsuffisantException {
        if (montant > e.solde) {
            throw new SoldeInsuffisantException();
        }
        e.solde -= montant;
        System.out.println("Nouveau solde : " + e.solde);
    }
}
\end{lstlisting}

\subsection{Explication}

\texttt{SoldeInsuffisantException} hérite de \texttt{Exception} (checked).
Le mot-clé \texttt{throws} dans la signature de \texttt{retirer()} oblige
l'appelant à gérer cette situation.

\section{Exercice 2 -- Montant Invalide}

\subsection{Objectif}

Créer une \texttt{MontantInvalideException} levée lorsque le montant est $\leq 0$.

\subsection{Code source}

\begin{lstlisting}[language=Java, caption={exo2\_2.java -- MontantInvalideException}]
class MontantInvalideException extends RuntimeException {}

public class exo2_2 {

    String code;
    double solde;

    public exo2_2(String code, double solde) {
        this.code = code;
        this.solde = solde;
    }

    void verser(exo2_2 e, double montant) {
        if (montant <= 0) {
            throw new MontantInvalideException();
        }
        e.solde += montant;
        System.out.println("Nouveau solde : " + e.solde);
    }
}
\end{lstlisting}

\subsection{Explication}

\texttt{MontantInvalideException} hérite de \texttt{RuntimeException}
(unchecked), car un montant invalide est une erreur de logique qui ne devrait
jamais se produire en production correcte.

\section{Exercice 3 -- Double Validation}

\subsection{Objectif}

Modifier \texttt{retirer()} pour gérer à la fois \texttt{MontantInvalideException}
et \texttt{SoldeInsuffisantException}.

\subsection{Code source}

\begin{lstlisting}[language=Java, caption={exo3\_2.java -- Double validation}]
class SoldeInsuffisantException3 extends Exception {}
class MontantInvalideException3 extends RuntimeException {}

public class exo3_2 {

    String code;
    double solde;

    public exo3_2(String code, double solde) {
        this.code = code;
        this.solde = solde;
    }

    void retirer(exo3_2 e, double montant) throws SoldeInsuffisantException3 {
        if (montant <= 0) {
            throw new MontantInvalideException3();
        }
        if (montant > e.solde) {
            throw new SoldeInsuffisantException3();
        }
        e.solde -= montant;
        System.out.println("Nouveau solde : " + e.solde);
    }
}
\end{lstlisting}

\section{Exercice 4 -- Inscription Utilisateur}

\subsection{Objectif}

Valider l'email et l'âge lors d'une inscription via deux exceptions personnalisées.

\subsection{Code source}

\begin{lstlisting}[language=Java, caption={exo4\_2.java -- EmailInvalideException \& AgeInvalideException}]
class EmailInvalideException4 extends Exception {}
class AgeInvalideException4 extends Exception {}

public class exo4_2 {

    void inscrire(exo4_2 e, String email, int age)
            throws EmailInvalideException4, AgeInvalideException4 {
        if (!email.contains("@")) {
            throw new EmailInvalideException4();
        }
        if (age < 18) {
            throw new AgeInvalideException4();
        }
        System.out.println("Inscription reussie");
    }
}
\end{lstlisting}

\section{Exercice 5 -- Authentification}

\subsection{Objectif}

Lever une \texttt{AuthentificationException} avec le message
\og Identifiants incorrects \fg{} en cas de mauvais login.

\subsection{Code source}

\begin{lstlisting}[language=Java, caption={exo5\_2.java -- AuthentificationException}]
class AuthentificationException5 extends Exception {}

public class exo5_2 {

    void login(exo5_2 e, String username, String password)
            throws AuthentificationException5 {
        if (!username.equals("admin") || !password.equals("1234")) {
            throw new AuthentificationException5();
        }
        System.out.println("Connexion reussie");
    }

    public static void main(String[] args) {
        exo5_2 e = new exo5_2();
        try {
            e.login(e, "user", "wrong");
        } catch (AuthentificationException5 ex) {
            System.out.println("Identifiants incorrects");
        }
    }
}
\end{lstlisting}

\section{Exercice 6 -- Gestion de Stock}

\subsection{Objectif}

Lever une \texttt{StockInsuffisantException} si la quantité demandée dépasse le stock disponible.

\subsection{Code source}

\begin{lstlisting}[language=Java, caption={exo6\_2.java -- StockInsuffisantException}]
class StockInsuffisantException6 extends Exception {}

public class exo6_2 {
    int stock;

    public exo6_2(int stock) {
        this.stock = stock;
    }

    void retirerDuStock(exo6_2 p, int quantite)
            throws StockInsuffisantException6 {
        if (quantite > p.stock) {
            throw new StockInsuffisantException6();
        }
        p.stock -= quantite;
        System.out.println("Stock restant : " + p.stock);
    }
}
\end{lstlisting}

\section{Exercice 7 -- Téléchargement}

\subsection{Objectif}

Lever une \texttt{QuotaDepasseException} si la taille d'un fichier dépasse la limite autorisée.

\subsection{Code source}

\begin{lstlisting}[language=Java, caption={exo7\_2.java -- QuotaDepasseException}]
class QuotaDepasseException7 extends Exception {}

public class exo7_2 {
    double limite = 500.0;

    void telechargerFichier(exo7_2 e, double taille)
            throws QuotaDepasseException7 {
        if (taille > e.limite) {
            throw new QuotaDepasseException7();
        }
        System.out.println("Telechargement fini");
    }
}
\end{lstlisting}

\section{Exercice 8 -- Validation de Formulaire}

\subsection{Objectif}

Lever une \texttt{ChampObligatoireException} si un champ du formulaire est vide ou nul.

\subsection{Code source}

\begin{lstlisting}[language=Java, caption={exo8\_2.java -- ChampObligatoireException}]
class ChampObligatoireException8 extends Exception {}

public class exo8_2 {

    void validerFormulaire(exo8_2 e, String nom, String email)
            throws ChampObligatoireException8 {
        if (nom == null || nom.isEmpty() || email == null || email.isEmpty()) {
            throw new ChampObligatoireException8();
        }
        System.out.println("Formulaire valide");
    }
}
\end{lstlisting}

\section{Exercice 9 -- Paiement}

\subsection{Objectif}

Gérer deux exceptions lors d'un paiement~: dépassement de plafond et carte expirée.

\subsection{Code source}

\begin{lstlisting}[language=Java, caption={exo9\_2.java -- PaiementRefuseException \& CarteExpireeException}]
class PaiementRefuseException9 extends Exception {}
class CarteExpireeException9 extends Exception {}

public class exo9_2 {
    double plafond = 2000.0;

    void payer(exo9_2 e, double montant, boolean expiree)
            throws PaiementRefuseException9, CarteExpireeException9 {
        if (expiree) {
            throw new CarteExpireeException9();
        }
        if (montant > e.plafond) {
            throw new PaiementRefuseException9();
        }
        System.out.println("Paiement effectue");
    }
}
\end{lstlisting}

\section{Exercice 10 -- Conception~: Quand Utiliser Chaque Type ?}

\subsection{Objectif}

Répondre aux questions de conception sur le choix du type d'exception.

\subsection{Synthèse}

\begin{table}[H]
\centering
\begin{tabularx}{0.95\textwidth}{L{4cm} X}
\toprule
\textbf{\color{primarycolor}Question} &
\textbf{\color{primarycolor}Réponse} \\
\midrule
Quand utiliser \textit{checked} ? &
Pour les erreurs \textbf{récupérables} et externes~: fichier introuvable,
connexion réseau, base de données. Le compilateur force la gestion. \\
\addlinespace[0.2cm]
Quand utiliser \textit{unchecked} ? &
Pour les \textbf{bugs logiques} du développeur~: index hors limites,
argument nul, montant invalide. Pas de \texttt{throws} requis. \\
\addlinespace[0.2cm]
Pourquoi créer une exception personnalisée ? &
Pour donner un \textbf{sens métier} précis à l'erreur, améliorer la
lisibilité du code et faciliter la maintenance. \\
\bottomrule
\end{tabularx}
\caption{Synthèse sur la conception des exceptions}
\end{table}

\chapter{Généricité}
\sloppy

\section{Exercice 1 -- Classe Générique \texttt{Boite<T>}}

\subsection{Objectif}

Créer une classe générique \texttt{Boite<T>} avec les méthodes \texttt{setContenu}
et \texttt{getContenu}, testée avec \texttt{String} et \texttt{Integer}.

\subsection{Code source}

\begin{lstlisting}[language=Java, caption={exo1\_3.java -- Boite générique}]
public class exo1_3<T> {
    T contenu;

    void setContenu(exo1_3<T> e, T contenu) {
        e.contenu = contenu;
    }

    T getContenu(exo1_3<T> e) {
        return e.contenu;
    }

    public static void main(String[] args) {
        exo1_3<String> boiteS = new exo1_3<>();
        boiteS.setContenu(boiteS, "Java");
        System.out.println("Contenu String : " + boiteS.getContenu(boiteS));

        exo1_3<Integer> boiteI = new exo1_3<>();
        boiteI.setContenu(boiteI, 42);
        System.out.println("Contenu Integer : " + boiteI.getContenu(boiteI));
    }
}
\end{lstlisting}

\subsection{Explication}

Le paramètre de type \texttt{T} est résolu à la compilation selon le type
utilisé lors de l'instanciation. Cela garantit la \textbf{sûreté de type}
(\textit{type-safety}) sans duplication de code.

\section{Exercice 2 -- Classe \texttt{Paire<T, U>}}

\subsection{Objectif}

Créer une classe générique à deux paramètres de type avec une méthode
\texttt{afficherPaire()}.

\subsection{Code source}

\begin{lstlisting}[language=Java, caption={exo2\_3.java -- Paire générique}]
public class exo2_3<T, U> {
    T premier;
    U second;

    public exo2_3(T premier, U second) {
        this.premier = premier;
        this.second = second;
    }

    void afficherPaire(exo2_3<T, U> p) {
        System.out.println("Paire -> Premier: " + p.premier
                + ", Second: " + p.second);
    }

    public static void main(String[] args) {
        exo2_3<String, Integer> paire = new exo2_3<>("Java", 2026);
        paire.afficherPaire(paire);
    }
}
\end{lstlisting}

\section{Exercice 3 -- Interface Générique \texttt{Calcul<T>}}

\subsection{Objectif}

Définir une interface générique \texttt{Calcul<T>} et l'implémenter pour
\texttt{Integer} et \texttt{Double}.

\subsection{Code source}

\begin{lstlisting}[language=Java, caption={exo3\_3.java -- Interface Calcul générique}]
interface Calcul<T> {
    T addition(T a, T b);
}

class CalculInt implements Calcul<Integer> {
    public Integer addition(Integer a, Integer b) {
        return a + b;
    }
}

class CalculDouble implements Calcul<Double> {
    public Double addition(Double a, Double b) {
        return a + b;
    }
}
\end{lstlisting}

\section{Exercice 4 -- Interface \texttt{Comparateur<T>}}

\subsection{Objectif}

Définir une interface \texttt{Comparateur<T>} et l'implémenter pour
\texttt{Integer} (par valeur) et \texttt{String} (par longueur).

\subsection{Code source}

\begin{lstlisting}[language=Java, caption={exo4\_3.java -- Interface Comparateur}]
interface Comparateur4<T> {
    int comparer(T a, T b);
}

class ComparateurInt4 implements Comparateur4<Integer> {
    public int comparer(Integer a, Integer b) {
        return a.compareTo(b);
    }
}

class ComparateurString4 implements Comparateur4<String> {
    public int comparer(String a, String b) {
        return Integer.compare(a.length(), b.length());
    }
}
\end{lstlisting}

\subsection{Explication}

La valeur retournée par \texttt{comparer()} suit la convention Java~:
négatif si $a < b$, zéro si $a = b$, positif si $a > b$.

\section{Exercice 5 -- Deuxième Implémentation de \texttt{Comparateur<T>}}

\subsection{Objectif}

Réimplémenter l'interface \texttt{Comparateur<T>} dans un fichier indépendant
pour consolider la compréhension.

\subsection{Code source}

\begin{lstlisting}[language=Java, caption={exo5\_3.java -- Comparateur5}]
interface Comparateur5<T> {
    int comparer(T a, T b);
}

class ComparateurInt5 implements Comparateur5<Integer> {
    public int comparer(Integer a, Integer b) {
        return a.compareTo(b);
    }
}

class ComparateurString5 implements Comparateur5<String> {
    public int comparer(String a, String b) {
        return Integer.compare(a.length(), b.length());
    }
}
\end{lstlisting}

\section{Exercice 6 -- Méthodes Génériques Statiques}

\subsection{Objectif}

Créer deux méthodes statiques génériques~: \texttt{afficherTableau(T[])}
et \texttt{getPremier(T[])}.

\subsection{Code source}

\begin{lstlisting}[language=Java, caption={exo6\_3.java -- Méthodes génériques statiques}]
public class exo6_3 {

    public static <T> void afficherTableau(T[] tableau) {
        for (T element : tableau) {
            System.out.print(element + " ");
        }
        System.out.println();
    }

    public static <T> T getPremier(T[] tableau) {
        if (tableau.length > 0) {
            return tableau[0];
        }
        return null;
    }

    public static void main(String[] args) {
        Integer[] nombres = {10, 20, 30, 40, 50};
        String[] mots = {"Java", "Data", "ENSAH"};

        afficherTableau(nombres);
        System.out.println("Premier : " + getPremier(nombres));
        afficherTableau(mots);
        System.out.println("Premier : " + getPremier(mots));
    }
}
\end{lstlisting}

\section{Exercice 7 -- Borne Supérieure \texttt{<T extends Number>}}

\subsection{Objectif}

Créer une méthode générique bornée qui calcule la somme de deux nombres
de n'importe quel type numérique.

\subsection{Code source}

\begin{lstlisting}[language=Java, caption={exo7\_3.java -- Borne supérieure}]
public class exo7_3 {

    public static <T extends Number> double somme(T a, T b) {
        return a.doubleValue() + b.doubleValue();
    }

    public static void main(String[] args) {
        System.out.println("Somme : " + somme(5, 3.14));
    }
}
\end{lstlisting}

\subsection{Explication}

\begin{accentbox}
    \centering\vspace{0.3cm}
    {\large La borne \texttt{<T extends Number>} restreint le paramètre
    de type aux sous-classes de \texttt{Number} (\texttt{Integer},
    \texttt{Double}, \texttt{Float}\ldots), permettant d'appeler
    \texttt{doubleValue()} en toute sécurité.}
    \vspace{0.3cm}
\end{accentbox}

\section{Exercice 8 -- Héritage Générique (Animal / Chien)}

\subsection{Objectif}

Créer une classe générique \texttt{Animal<T>} et une sous-classe
\texttt{Chien} qui fixe le type à \texttt{String}.

\subsection{Code source}

\begin{lstlisting}[language=Java, caption={exo8\_3.java -- Héritage générique}]
class Animal<T> {
    T nom;
}

class Chien extends Animal<String> {
    public Chien(String nom) {
        this.nom = nom;
    }
}

public class exo8_3 {
    public static void main(String[] args) {
        Chien c = new Chien("Rex");
        System.out.println("Nom du chien : " + c.nom);
    }
}
\end{lstlisting}

\section{Exercice 9 -- Héritage Paramétré (Vehicule / Voiture)}

\subsection{Objectif}

Créer une classe générique \texttt{Vehicule<T>} et une sous-classe
\texttt{Voiture<T>} qui hérite du paramètre de type.

\subsection{Code source}

\begin{lstlisting}[language=Java, caption={exo9\_3.java -- Vehicule et Voiture génériques}]
class Vehicule<T> {
    T vitesse;
}

class Voiture<T> extends Vehicule<T> {
    public Voiture(T vitesse) {
        this.vitesse = vitesse;
    }
}

public class exo9_3 {
    public static void main(String[] args) {
        Voiture<Integer> v = new Voiture<>(120);
        System.out.println("Vitesse : " + v.vitesse);
    }
}
\end{lstlisting}

\section{Exercice 10 -- Repository Générique}

\subsection{Objectif}

Créer un \texttt{Repository<T>} générique et une classe
\texttt{UserRepository} qui le spécialise pour les objets \texttt{User}.

\subsection{Code source}

\begin{lstlisting}[language=Java, caption={exo10\_3.java -- Repository générique}]
class User {
    String nom;
    public User(String nom) { this.nom = nom; }
    public String toString() { return "User: " + nom; }
}

class Repository<T> {
    void save(T obj) {
        System.out.println("Objet sauvegarde : " + obj.toString());
    }
}

class UserRepository extends Repository<User> {}

public class exo10_3 {
    public static void main(String[] args) {
        User u = new User("Yazid");
        UserRepository repo = new UserRepository();
        repo.save(u);
    }
}
\end{lstlisting}

\section{Exercice 11 -- Wildcard \texttt{<?>}}

\subsection{Objectif}

Créer une méthode \texttt{afficherListe(List<?> liste)} acceptant
n'importe quel type de liste.

\subsection{Code source}

\begin{lstlisting}[language=Java, caption={exo11\_3.java -- Wildcard ?}]
import java.util.List;
import java.util.ArrayList;

public class exo11_3 {

    void afficherListe(List<?> liste) {
        for (Object obj : liste) {
            System.out.print(obj + " ");
        }
        System.out.println();
    }

    public static void main(String[] args) {
        List<String> noms = new ArrayList<>();
        noms.add("Yazid");
        noms.add("ENSAH");

        exo11_3 e = new exo11_3();
        e.afficherListe(noms);
    }
}
\end{lstlisting}

\subsection{Explication}

Le wildcard \texttt{<?>} signifie \og n'importe quel type \fg{}. Il permet
d'écrire une méthode générique sans spécifier de paramètre de type, mais
en contrepartie on ne peut \textbf{pas} ajouter d'éléments à la liste
(sauf \texttt{null}).

\section{Exercice 12 -- Wildcard Borné \texttt{<? extends Number>}}

\subsection{Objectif}

Créer une méthode \texttt{afficherNombres(List<? extends Number>)} testée
avec \texttt{List<Integer>} et \texttt{List<Double>}.

\subsection{Code source}

\begin{lstlisting}[language=Java, caption={exo12\_3.java -- Wildcard borné}]
import java.util.List;
import java.util.ArrayList;

public class exo12_3 {

    void afficherNombres(List<? extends Number> liste) {
        for (Number n : liste) {
            System.out.print(n + " ");
        }
        System.out.println();
    }

    public static void main(String[] args) {
        List<Integer> entiers = new ArrayList<>();
        entiers.add(10);
        entiers.add(20);

        List<Double> reels = new ArrayList<>();
        reels.add(15.5);
        reels.add(25.5);

        exo12_3 e = new exo12_3();
        e.afficherNombres(entiers);
        e.afficherNombres(reels);
    }
}
\end{lstlisting}

\subsection{Comparaison des Wildcards}

\begin{table}[H]
\centering
\begin{tabularx}{0.95\textwidth}{C{3.5cm} X X}
\toprule
\textbf{\color{primarycolor}Wildcard} &
\textbf{\color{primarycolor}Signification} &
\textbf{\color{primarycolor}Usage typique} \\
\midrule
\texttt{<?>} & N'importe quel type & Lecture seule, type inconnu \\
\addlinespace[0.15cm]
\texttt{<? extends T>} & T ou ses sous-types & Lecture d'une collection de T \\
\addlinespace[0.15cm]
\texttt{<? super T>} & T ou ses super-types & Écriture dans une collection \\
\bottomrule
\end{tabularx}
\caption{Comparaison des wildcards en Java}
\end{table}


% ── Conclusion ─────────────────────────────────────────────────────────────
\chapter*{Conclusion}
\addcontentsline{toc}{chapter}{Conclusion}

Ce travail pratique nous a permis de maîtriser deux piliers fondamentaux de la programmation orientée objet en Java : la \textbf{gestion des exceptions} et la \textbf{généricité}.

La première partie a mis en évidence la différence entre exceptions vérifiées (\textit{checked}) et non vérifiées (\textit{unchecked}), ainsi que l'importance de créer des hiérarchies d'exceptions métier claires et expressives. La propagation d'exceptions à travers la pile d'appels et l'usage du mot-clé \texttt{throw} ont été mis en pratique de manière concrète.

La deuxième partie a montré la puissance des exceptions personnalisées pour rendre le code plus lisible, robuste et maintenable, en donnant un sens métier précis aux erreurs détectées à l'exécution.

Enfin, la troisième partie a illustré l'intérêt de la généricité pour écrire du code réutilisable, type-safe, applicable à des structures de données, des interfaces et des classes génériques avec bornes et wildcards.

\vspace{2cm}
\begin{flushright}
\textit{\docAuthor}\\
\textit{ENSAH -- \docYear}
\end{flushright}

\end{document}
